% GENERATED FILE. Edit canonical lessons, then run npm run generate:books.
% canonical-source-hash: fnv1a64:219b2452f2696e25
% canonical-lessons: BN-C489-bagdan, BN-C489-songsar, BN-C489-man-bhangano, BN-C489-ashirbad, BN-C489-pronam-kora

\chapter{Weddings and Households}
\label{ch:bn-weddings-and-households}
\hlchaptermodality{489}

\begin{chapteropening}
\textbf{By the end of this chapter:} \emph{I can talk about an engagement, a household and a blessing.}

Complete the last lesson of chapter 489: I can talk about an engagement, a household and a blessing.
\end{chapteropening}

\section[\texorpdfstring{\bn{বাগদান} (bāgdān)}{বাগদান (bāgdān)}]{\bn{বাগদান} (bāgdān) — an engagement (to marry)}

\label{lesson:BN-C489-bagdan}



\begin{quote}
Before the new one: say the Bengali for old age, then the Bengali for to grow up.
\end{quote}

\subsection*{You'll want to know: \bn{বাগদান}}
\textbf{\bn{বাগদান}} — \emph{bāgdān} — \textquotedblleft{}an engagement (to marry)\textquotedblright{}.

\begin{quote}
\bn{গত} \bn{মাসে} \bn{মীরার} \bn{বাগদান} \bn{হল।}
\end{quote}

\begin{quote}
\emph{gôtô māse mīrār bāgdān hôlô.}
\end{quote}

\begin{quote}
Mira got engaged last month.
\end{quote}

\begin{cousinweb}
\textbf{\bn{বাগদান}} (\emph{bāgdān}) is a tatsama, Sanskrit \emph{vāgdāna}, \textquotedblleft{}a gift by word of mouth\textquotedblright{}, the promise of a bride in marriage (Monier-Williams).
\end{cousinweb}

\begin{grammarlens}[title={What you've built}]
One more word for describing what you lived and felt.
\end{grammarlens}

\subsection*{Your turn}
\begin{itemize}
\raggedright
  \item \emph{Say it:} \emph{bāgdān}
  \item \emph{Say it:} \emph{bāgdān}, once more
  \item \emph{Say it:} \emph{gôtô māse mīrār bāgdān hôlô.} — Mira got engaged last month.
  \item \emph{Recall:} say the Bengali for old age, then the Bengali for to grow up, then say \emph{bāgdān} again
\end{itemize}

\begin{tcolorbox}[breakable,colback=teal!4,colframe=teal!35!black,title={Before you move on}]
What does \bn{বাগদান} mean? (\textquotedblleft{}An engagement (to marry)\textquotedblright{}.) Say it once more.
\end{tcolorbox}
\section[\texorpdfstring{\bn{সংসার} (sôngsār)}{সংসার (sôngsār)}]{\bn{সংসার} (sôngsār) — a household; family life}

\label{lesson:BN-C489-songsar}



\begin{quote}
Before the new one: say the Bengali for to grow up, then the Bengali for an engagement.
\end{quote}

\subsection*{You'll want to know: \bn{সংসার}}
\textbf{\bn{সংসার}} — \emph{sôngsār} — \textquotedblleft{}a household; family life\textquotedblright{}.

\begin{quote}
\bn{বিয়ের} \bn{পরে} \bn{মীরা} \bn{নতুন} \bn{সংসার} \bn{শুরু} \bn{করল।}
\end{quote}

\begin{quote}
\emph{biyer pôre mīrā nôtun sôngsār shuru kôrlô.}
\end{quote}

\begin{quote}
After the wedding Mira set up a new household.
\end{quote}

\begin{cousinweb}
\textbf{\bn{সংসার}} (\emph{sôngsār}) is a tatsama, Sanskrit \emph{saṃsāra}, \textquotedblleft{}wandering through\textquotedblright{}, the round of worldly life (Monier-Williams); Bengali brought it down to the household and everyday family life.
\end{cousinweb}

\begin{grammarlens}[title={What you've built}]
One more word for describing what you lived and felt.
\end{grammarlens}

\subsection*{Your turn}
\begin{itemize}
\raggedright
  \item \emph{Say it:} \emph{sôngsār}
  \item \emph{Say it:} \emph{sôngsār}, once more
  \item \emph{Say it:} \emph{biyer pôre mīrā nôtun sôngsār shuru kôrlô.} — After the wedding Mira set up a new household.
  \item \emph{Recall:} say the Bengali for to grow up, then the Bengali for an engagement, then say \emph{sôngsār} again
\end{itemize}

\begin{tcolorbox}[breakable,colback=teal!4,colframe=teal!35!black,title={Before you move on}]
What does \bn{সংসার} mean? (\textquotedblleft{}A household; family life\textquotedblright{}.) Say it once more.
\end{tcolorbox}
\section[\texorpdfstring{\bn{মান} \bn{ভাঙানো} (mān bhāngāno)}{মান ভাঙানো (mān bhāngāno)}]{\bn{মান} \bn{ভাঙানো} (mān bhāngāno) — to coax someone out of a sulk}

\label{lesson:BN-C489-man-bhangano}



\begin{quote}
Before the new one: say the Bengali for an engagement, then the Bengali for a household.
\end{quote}

\subsection*{You'll want to know: \bn{মান} \bn{ভাঙানো}}
\textbf{\bn{মান} \bn{ভাঙানো}} — \emph{mān bhāngāno} — \textquotedblleft{}to coax someone out of a sulk\textquotedblright{}.

\begin{quote}
\bn{স্বামী} \bn{ফুল} \bn{এনে} \bn{স্ত্রীর} \bn{মান} \bn{ভাঙালেন।}
\end{quote}

\begin{quote}
\emph{svāmī phul ene strīr mān bhāngālen.}
\end{quote}

\begin{quote}
The husband brought flowers and coaxed his wife out of her sulk.
\end{quote}

Past (he/she …): \bn{তিনি} \bn{ছেলের} \bn{মান} \bn{ভাঙালেন}, \emph{tini chheler mān bhāngālen}.

\begin{cousinweb}
\textbf{\bn{মান} \bn{ভাঙানো}} (\emph{mān bhāngāno}) is \emph{mān}, wounded pride, which CDIAL lists for Bengali under Sanskrit \emph{māna} (CDIAL 10040), with \emph{bhāngāno}, \textquotedblleft{}to break\textquotedblright{}, the causative of \emph{bhāngā} (CDIAL 9353): to break the sulk.
\end{cousinweb}

\begin{grammarlens}[title={What you've built}]
One more word for describing what you lived and felt.
\end{grammarlens}

\subsection*{Your turn}
\begin{itemize}
\raggedright
  \item \emph{Say it:} \emph{mān bhāngāno}
  \item \emph{Say it:} \emph{mān bhāngāno}, once more
  \item \emph{Say it:} \emph{svāmī phul ene strīr mān bhāngālen.} — The husband brought flowers and coaxed his wife out of her sulk.
  \item \emph{Recall:} say the Bengali for an engagement, then the Bengali for a household, then say \emph{mān bhāngāno} again
\end{itemize}

\begin{tcolorbox}[breakable,colback=teal!4,colframe=teal!35!black,title={Before you move on}]
What does \bn{মান} \bn{ভাঙানো} mean? (\textquotedblleft{}To coax someone out of a sulk\textquotedblright{}.) Say it once more.
\end{tcolorbox}
\section[\texorpdfstring{\bn{আশীর্বাদ} (āshīrbād)}{আশীর্বাদ (āshīrbād)}]{\bn{আশীর্বাদ} (āshīrbād) — a blessing}

\label{lesson:BN-C489-ashirbad}



\begin{quote}
Before the new one: say the Bengali for a household, then the Bengali for to coax someone out of a sulk.
\end{quote}

\subsection*{You'll want to know: \bn{আশীর্বাদ}}
\textbf{\bn{আশীর্বাদ}} — \emph{āshīrbād} — \textquotedblleft{}a blessing\textquotedblright{}.

\begin{quote}
\bn{বিয়ের} \bn{দিন} \bn{দাদু} \bn{বর} \bn{আর} \bn{কনেকে} \bn{আশীর্বাদ} \bn{করলেন।}
\end{quote}

\begin{quote}
\emph{biyer din dādu bôr ār kôneke āshīrbād kôrlen.}
\end{quote}

\begin{quote}
On the wedding day my grandfather blessed the bride and groom.
\end{quote}

\bn{আশীর্বাদ} \bn{করা}, to bless.

\begin{cousinweb}
\textbf{\bn{আশীর্বাদ}} (\emph{āshīrbād}) is a tatsama, Sanskrit \emph{āśīrvāda}, \textquotedblleft{}benediction\textquotedblright{} (Monier-Williams): the saying of a wish.
\end{cousinweb}

\begin{grammarlens}[title={What you've built}]
One more word for describing what you lived and felt.
\end{grammarlens}

\subsection*{Your turn}
\begin{itemize}
\raggedright
  \item \emph{Say it:} \emph{āshīrbād}
  \item \emph{Say it:} \emph{āshīrbād}, once more
  \item \emph{Say it:} \emph{biyer din dādu bôr ār kôneke āshīrbād kôrlen.} — On the wedding day my grandfather blessed the bride and groom.
  \item \emph{Recall:} say the Bengali for a household, then the Bengali for to coax someone out of a sulk, then say \emph{āshīrbād} again
\end{itemize}

\begin{tcolorbox}[breakable,colback=teal!4,colframe=teal!35!black,title={Before you move on}]
What does \bn{আশীর্বাদ} mean? (\textquotedblleft{}A blessing\textquotedblright{}.) Say it once more.
\end{tcolorbox}
\section[\texorpdfstring{\bn{প্রণাম} \bn{করা} (prônām kôrā)}{প্রণাম করা (prônām kôrā)}]{\bn{প্রণাম} \bn{করা} (prônām kôrā) — to touch an elder's feet; to bow in respect}

\label{lesson:BN-C489-pronam-kora}



\begin{quote}
Before the new one: say the Bengali for to coax someone out of a sulk, then the Bengali for a blessing.
\end{quote}

\subsection*{You'll want to know: \bn{প্রণাম} \bn{করা}}
\textbf{\bn{প্রণাম} \bn{করা}} — \emph{prônām kôrā} — \textquotedblleft{}to touch an elder's feet; to bow in respect\textquotedblright{}.

\begin{quote}
\bn{নতুন} \bn{বছরে} \bn{অরুণ} \bn{দাদুকে} \bn{প্রণাম} \bn{করল।}
\end{quote}

\begin{quote}
\emph{nôtun bôchhôre ôrun dāduke prônām kôrlô.}
\end{quote}

\begin{quote}
At the start of the new year, Arun touched his grandfather's feet.
\end{quote}

Past (he/she …): \bn{তিনি} \bn{ঠাকুমাকে} \bn{প্রণাম} \bn{করলেন}, \emph{tini ṭhākumāke prônām kôrlen}.

\begin{cousinweb}
\textbf{\bn{প্রণাম} \bn{করা}} (\emph{prônām kôrā}) is Sanskrit \emph{praṇāma}, \textquotedblleft{}bowing, a respectful salutation\textquotedblright{} (Monier-Williams), with \emph{kôrā}, \textquotedblleft{}to do\textquotedblright{} (CDIAL 2814).
\end{cousinweb}

\begin{grammarlens}[title={What you've built}]
One more word for describing what you lived and felt.
\end{grammarlens}

\subsection*{Your turn}
\begin{itemize}
\raggedright
  \item \emph{Say it:} \emph{prônām kôrā}
  \item \emph{Say it:} \emph{prônām kôrā}, once more
  \item \emph{Say it:} \emph{nôtun bôchhôre ôrun dāduke prônām kôrlô.} — At the start of the new year, Arun touched his grandfather's feet.
  \item \emph{Recall:} say the Bengali for to coax someone out of a sulk, then the Bengali for a blessing, then say \emph{prônām kôrā} again
\end{itemize}

\begin{tcolorbox}[breakable,colback=teal!4,colframe=teal!35!black,title={Before you move on}]
What does \bn{প্রণাম} \bn{করা} mean? (\textquotedblleft{}To touch an elder's feet; to bow in respect\textquotedblright{}.) Say it once more.
\end{tcolorbox}
